% BEGIN VOL07-EXPANSION VII18
\section{Supplementary worked examples}
\label{sec:vii18-pedagogy}

\begin{example}[Cylinder as a developable ruled surface]\label{ex:vii18-ped-01}
The cylinder
\[
X(u,v)=c(u)+v\,d,\qquad c(u)=(R\cos u,R\sin u,0),\quad d=(0,0,1),
\]
is ruled by parallel lines.  One principal curvature vanishes, so \(K=0\).  A sheet can be rolled into a cylinder without stretching, reflecting local developability.
\end{example}

\begin{example}[Cone away from the apex]\label{ex:vii18-ped-02}
A cone can be parametrized by \(X(u,v)=v(\cos u,\sin u,a)\) for \(v>0\).  It is ruled and has \(K=0\) away from the apex, but the apex is singular.  Developability is therefore a local smooth statement that does not automatically include singular points.
\end{example}

\begin{example}[Helicoid: ruled but not developable]\label{ex:vii18-ped-03}
The helicoid
\[
X(u,v)=(v\cos u,v\sin u,au)
\]
is ruled because \(v\mapsto X(u,v)\) is a line for fixed \(u\).  Its Gaussian curvature is negative away from infinity, so ruled does not imply developable.
\end{example}

\section{Graded supplementary exercises}
The following exercises add a deliberate progression from concrete calculation to structural proof, hypothesis testing, application, and synthesis.

\subsection*{Standard computations and constructions}

\begin{exercise}[Cylinder rulings]\label{exr:vii18-ped-01}
Identify the ruling direction in \(X(u,v)=(R\cos u,R\sin u,v)\).
\end{exercise}
\begin{hint}
Fix \(u\) and vary \(v\).
\end{hint}
\begin{solution}
The rulings are vertical lines with direction \((0,0,1)\).
\end{solution}

\begin{exercise}[Cone rulings]\label{exr:vii18-ped-02}
For \(X(u,v)=v(\cos u,\sin u,a)\), describe the ruling through fixed \(u\).
\end{exercise}
\begin{hint}
Vary \(v\).
\end{hint}
\begin{solution}
It is the straight ray through the origin in direction \((\cos u,\sin u,a)\).
\end{solution}

\begin{exercise}[Helicoid rulings]\label{exr:vii18-ped-03}
Show \(v\mapsto(v\cos u,v\sin u,au)\) is a straight line for fixed \(u\).
\end{exercise}
\begin{hint}
Separate constant and linear terms in \(v\).
\end{hint}
\begin{solution}
It is \((0,0,au)+v(\cos u,\sin u,0)\), hence a line.
\end{solution}

\begin{exercise}[Developable curvature]\label{exr:vii18-ped-04}
What Gaussian curvature does a smooth developable surface have?
\end{exercise}
\begin{hint}
Use the local isometry to the plane or rank-one shape operator.
\end{hint}
\begin{solution}
It has \(K=0\) at regular points.
\end{solution}

\begin{exercise}[Cylinder unrolling]\label{exr:vii18-ped-05}
Give local plane coordinates that preserve the metric of a radius-\(R\) cylinder.
\end{exercise}
\begin{hint}
Use \(x=Ru,\ y=v\).
\end{hint}
\begin{solution}
The cylinder metric \(R^2du^2+dv^2\) becomes \(dx^2+dy^2\).
\end{solution}

\subsection*{Proofs}

\begin{exercise}[Cylinder has zero Gaussian curvature]\label{exr:vii18-ped-06}
Prove this from its principal directions.
\end{exercise}
\begin{hint}
One principal curvature is zero along rulings.
\end{hint}
\begin{solution}
The axial normal does not change, so one principal curvature is \(0\); therefore \(K=k_1k_2=0\).
\end{solution}

\begin{exercise}[Developable local isometry implies \(K=0\)]\label{exr:vii18-ped-07}
Use the intrinsic invariance of Gaussian curvature.
\end{exercise}
\begin{hint}
A plane has zero Gaussian curvature.
\end{hint}
\begin{solution}
By Gauss' Theorema Egregium, local isometries preserve \(K\); a surface locally isometric to the plane must therefore have \(K=0\).
\end{solution}

\begin{exercise}[Tangent plane criterion along rulings]\label{exr:vii18-ped-08}
Explain why a regular ruled surface is developable when its tangent plane is constant along each ruling.
\end{exercise}
\begin{hint}
Translate constant tangent plane into constant normal along each ruling.
\end{hint}
\begin{solution}
The normal derivative vanishes in the ruling direction, so the shape operator has a zero eigenvalue there and hence \(K=0\); under standard regularity hypotheses this is the developable condition.
\end{solution}

\begin{exercise}[Cone singularity]\label{exr:vii18-ped-09}
Show the cone parametrization loses rank at the apex \(v=0\).
\end{exercise}
\begin{hint}
Compute \(X_u\) and \(X_v\).
\end{hint}
\begin{solution}
At \(v=0\), \(X_u=0\), so the derivative cannot have rank two.
\end{solution}

\subsection*{Counterexamples and hypothesis tests}

\begin{exercise}[Every ruled surface is developable]\label{exr:vii18-ped-10}
Test with the helicoid.
\end{exercise}
\begin{hint}
Recall its negative Gaussian curvature.
\end{hint}
\begin{solution}
False.  The helicoid is ruled but has \(K<0\) at regular points.
\end{solution}

\begin{exercise}[Every \(K=0\) surface is globally a plane]\label{exr:vii18-ped-11}
Test with a cylinder.
\end{exercise}
\begin{hint}
Local flatness does not imply global embedding as a plane.
\end{hint}
\begin{solution}
False.  A cylinder has \(K=0\) but is globally wrapped.
\end{solution}

\begin{exercise}[Cone apex is a regular developable point]\label{exr:vii18-ped-12}
Test the rank at the apex.
\end{exercise}
\begin{hint}
The parametrization collapses there.
\end{hint}
\begin{solution}
False.  The apex is singular and is excluded from the regular surface.
\end{solution}

\subsection*{Applications and investigations}

\begin{exercise}[Sheet-metal bending]\label{exr:vii18-ped-13}
Why are developable surfaces important in fabrication from flat sheets?
\end{exercise}
\begin{hint}
Local isometry preserves lengths in the sheet.
\end{hint}
\begin{solution}
They can be formed by bending without first-order stretching, so planar material patterns transfer naturally to cylinders and cones.
\end{solution}

\begin{exercise}[Texture mapping]\label{exr:vii18-ped-14}
Why can a developable patch be flattened with little intrinsic distortion?
\end{exercise}
\begin{hint}
Its metric is locally Euclidean.
\end{hint}
\begin{solution}
Local isometry to the plane preserves intrinsic lengths and angles, making planar parameterization especially natural.
\end{solution}

\subsection*{Challenge problems}

\begin{exercise}[Tangent developable]\label{exr:vii18-ped-15}
For a regular space curve \(\gamma(s)\), consider \(X(s,v)=\gamma(s)+vT(s)\).  Explain why this surface is ruled and where singularities can occur.
\end{exercise}
\begin{hint}
Differentiate with respect to \(s\) and \(v\).
\end{hint}
\begin{solution}
Rulings are tangent lines.  Since \(X_v=T\) and \(X_s=T+v\kappa N\), the two derivatives are dependent at \(v=0\), so the original curve is a singular locus of the tangent developable when \(\kappa\neq0\).
\end{solution}

\begin{exercise}[Why the helicoid is not flattenable]\label{exr:vii18-ped-16}
Use Gaussian curvature to explain why a helicoid cannot be locally isometric to the plane at a point with \(K<0\).
\end{exercise}
\begin{hint}
Invoke intrinsic invariance of \(K\).
\end{hint}
\begin{solution}
A plane has \(K=0\), and local isometries preserve Gaussian curvature, so negative-curvature helicoid points cannot be flattened isometrically.
\end{solution}

% END VOL07-EXPANSION VII18
